\section{Kết luận và hướng phát
triển}\label{sec:extra0}

Bài báo đã trình bày một đánh giá bán thực nghiệm đối với cơ chế
Hash-linked Tamper-evident Ledger nhằm kiểm chứng tính toàn vẹn dữ liệu
trong hệ thống Data Pipeline. Cơ chế này ghi nhận bằng chứng integrity ở
từng giai đoạn và liên kết các ledger blocks thông qua previous hash.
Verification engine tính toán lại số lượng bản ghi, mã băm cấu trúc dữ liệu (schema hash), batch hash,
block hash và liên kết chuỗi để phát hiện modification sau xử lý.

Các quan sát thực nghiệm hỗ trợ các giả thuyết chính trong phạm vi
bounded scope. Tất cả tamper scenarios được định nghĩa trước đều được
phát hiện, điều kiện sạch vẫn VALID trong bảng điều khiển xuất dữ liệu quan sát
được, và lineage context hỗ trợ xác định affected giai đoạn. Runtime
chi phí phát sinh tăng khi số lượng bản ghi tăng, nhưng kết quả chi phí phát sinh vẫn là
approximate và chịu ảnh hưởng bởi cloud compute variability.

Nghiên cứu nên được hiểu là bằng chứng cho một cơ chế tamper-evident
nhẹ, không phải bằng chứng về bảo mật Blockchain đầy đủ. Nguyên mẫu
không bao gồm decentralized consensus, nhiều independent nút, smart
contracts hoặc external neo tin cậy. Vì vậy, security boundary của
cơ chế hiện tại bị giới hạn trong controlled không gian làm việc setting.

Trong tương lai, nghiên cứu cần được lặp lại trên nhiều data platforms,
mở rộng tập tamper vectors thông qua kiểm thử đối kháng, tăng số lần
số lần lặp lại đo hiệu năng, xuất dữ liệu raw benchmark values, externalize ledger
sang một ranh giới tin cậy độc lập và bổ sung chữ ký số nhằm tăng
cường bằng chứng chống lại privileged người có quyền truy cập nội bộ modification.
